\section{The signed short arcs at inert primes}

% BEGIN U83M1
\paragraph{U83M1. Simultaneous Wach lifts before every nonthick letter.}
Fix an actual mixed C or D orbit with the word of I82W2--I82W5.
Write $d=n_w$, $q=p-1$, $w=W_0$, $S=S_0$, $\eta=H$,
$h=B+\eta+S$, $G=w+S$, and put
\[
 D_0=G+\eta,\quad L_0=\log_p(2h),\quad
 U=h/p^2+D_0/p+L_0+1 .
\]
These are the mixed branch of the native height and its $\mathrm{Sg}$.
In $\mathrm{arcOf}$ every letter retains its sign and size, even
when its host position is nonthick. Cut just before every nonthick
position. M82R1 gives lengths $1\le e\le3$ summing to $d$.
The first letter after a cut has size at most $w$, and $s_i>\eta$
everywhere. Thus each signed window holds for the source and target
and also for any crystalline extension of the target by the trivial
representation. The added zero lies in its central window. Its
underlying lattice projection onto $O$ is surjective.

At every cut the series of R82W3 for either lattice map starts
with height at most $D_0$ and all later heights are at most $h+\eta$.
Write $H_* =\max_i (M_i+h_i)$ for the input widths and
$c_H=\lceil1+\log_p(H_*+1)\rceil\le2+L_0$,
since $\eta+1\le h$. Summing the geometric series in the remaining
terms, R82W4 kills the invariant Wach cokernel at that cut at an
integer exponent at most
\[
 P\{D_0+c_H+3+(h+\eta+c_H+3)/(p-2)\}+c_H,
 \qquad P=p/q^2.
\]
This quantity is at most $2D_0/p+2h/p^2+2L_0+3\le3U$.
Indeed $P(1+1/(p-2))\le2/p$ and also $\le1/8$,
$P/(p-2)\le2/p^2$, and $\eta\le D_0$ since $p\ge11$.
The remaining constants are bounded by
$c_H+(c_H+3)/8\le 3+2L_0$.
One may therefore use the same $R=\lfloor3U\rfloor$ at all
cuts, for the tested surjection and for the trivial projection of
\emph{every} such target extension. The bound is uniform in its class.

The native starting precision and its truncation cost are
\[
 M=\lfloor 2h/q\rfloor+1,\qquad
 \sigma=\lceil(2G+M)P+1\rceil\le3U .
\]
For the last assertion use $2P\le3/p$, $2P/q\le3/p^2$,
$P\le1$ and $M\le2h/q+1$. Finally $40U\le B/4$ uniformly
on the late family: $h/B\to1$, $D_0/B\to0$, $p>x^{\epsilon_z/5}$,
and $\log(2h)/B\to0$ by I82W1--I82W3. This verifies the large-size
guard for the two possible length-three orientations, even if the
actual $p$ varies with the member.
% END U83M1

% BEGIN U83M2
\paragraph{U83M2. Integral decomposition with the signed-arc precision.}
Work with $A=O[[v]]$ and $P_s=\varphi^s(v+p)$ as in I82B4,
and write $L_i=N_i^\Delta$. The first forward step after a cut
uses the modulus $f=v^M P_0^{2G}$. For $z\in vL_c$ there are
$x\in vL_c$, $b\ge0$ and $\tau\in v A\varphi L_c$ such that
\[
 (\gamma-1)x-p^b z\in v^M L_c,\qquad
 p^b\tau-p^\sigma\varphi x\in fA\varphi L_c.
\]
To see the stated $v$ factor as well as the exact exponent, solve
for a primitive of $z$ with zero constant up to degree $M-1$.
Its degree-$k$ denominator divides $p^{a_k}$,
$a_k=k+v_p(k!)\le kp/q$, by the recursion of I82B4.
Put $D=2G+M$. The image of $v^k$ for $k\ge1$ is a unit times
$v^kP_0^{qk}$. If $qk\le D$, then $a_k\le\sigma$. Otherwise
remove $P_0^{2G}$ and expand $P_0^{qk-2G}=(p+v)^{qk-2G}$.
Modulo $v^M$ the expansion is divisible by $p^{qk-D}$.
Since $q>p/q$, $a_k\le\sigma+qk-D$. Thus in the quotient by $f$
every transformed term times $p^\sigma$ has an integral representative
with zero constant: retain the outside factor $v^k$ and perform the
expansion modulo $v^M$ before multiplying by it. Finite summation
constructs $\tau$. Take $x$ to be the polynomial primitive times a
common $p^b$; cancellation of any additional denominator modulo $f$
is valid by distinguished division. Commuting the two actions then gives
$p^\sigma\varphi z-(\gamma-1)\tau\in f A\varphi L_c$,
since $qM>2h\ge2G$ and $\gamma$ preserves both factor ideals.
The latter module lies in $v^M L_{c+1}$ by the height bound $G$
at the first forward letter. The recursion for the primitive to
all higher degrees exists over the fraction field and supplies the
solution in the native $\mathrm{IsTr}$, which only tests its first $M$ terms.

For an arc ending at $c$ of length $e$, put
$H_s= h$ for $s< e-1$, $H_{e-1}=G$, and
$\Lambda=\prod_{s<e}P_s^{H_s}$. These are the exponents of
$\mathrm{abelian}_2$, since every letter has signed type.
At $c-s$ the true upper and lower heights are both at most $H_s$.
Hence iterated height inclusions give
$\Lambda L_c\subset A\varphi^e L_{c-e}$ and
$\Lambda A\varphi^e L_{c-e}\subset L_c$, also in a target extension.
Use the cut lift of $p^R$ from U83M1 and the truncations above in
I82B5's telescoping identity with $j=1$. Its contraction is valid
since $qM>2h\ge h+1$ and each Frobenius denominator is at most
$P_0^h$. Its kernel identity uses the same flat Frobenius base
change as I82B5. This decomposes $p^{R+2\sigma}$ times every
integral finite target class as a sum of classes at these cuts.
Their cut vector $w_c$ satisfies $(\gamma-1)w_c\in vL_c$ and
$\Lambda w_c\in A\varphi^e L_{c-e}$. Each class is finite by
H82A1 since the second Herr coordinate lies in $v^ML$ and a
finite product of the $P_s$ clears the first. In particular the
preliminary decomposition costs at most $9U$, with the nonthick
height in this cost divided by $p$.
% END U83M2

% BEGIN U83M3
\paragraph{U83M3. The integral exceptional directions on all primitive levels.}
On this arc let $U_O=L_{c-e}/vL_{c-e}$ and transport its rational
constant space to the crystalline component at $c$ by $\varphi^e$.
Write $F_s^a$ for the transported de Rham filtration from $c-s$
there, as in D82P1. The map $U_O[1/p]\to D_c$ is invertible by
W79H3 and the isomorphism given by crystalline Frobenius. For each
up letter let $C_s=F_s^{s_{c-s}}$ and for each down letter let
$B_s=F_s^{-\eta}$. Then $\mathrm{codim}\,C_s\le m$ and
$\dim B_s\le m$ since the central window contains at most $m$
jumps and the cloud is disjoint from it. These assertions use the
small multiplicity at every letter, including the first nonthick one.
Split off $U_O\cap\sum B_s$ and in the complementary lattice split
off its intersection with $\bigcap C_s$, as in V81P1. Both quotients
are torsion-free. The resulting $U_O=U'\oplus S_r$ has
$U'_L\subset C_s$ at every up letter and $B_s\subset (S_r)_L$
at every down letter, with $k=\mathrm{rank}_O S_r\le me$.
A basis of this splitting lifts to a Wach basis, and its horizontal
series are the unique normalized solutions at $v=0$ from J82F1.

In these constant coordinates set $\Xi=\Lambda^{-1}A\otimes U_O$.
Let $\mathcal N\subset\Xi$ impose the conditions that all degree-$i$
primitive Taylor coefficients at level $s$ lie in $F_s^{-i}$,
with the coefficient algebra of its primitive roots understood.
This is the native $\mathrm{InNL}$ module. Indeed a native section
is a scalar times a constant vector $b$ with prescribed orders
$\alpha_s$ at all levels and $b\in F_s^{-\alpha_s}$, so satisfies
these tests. Conversely work modulo the ideal $(\Lambda^2)$
on the numerator after inverting $p$. The $P_s$ are pairwise
coprime and separable there, and distinguished division identifies
the quotient with the product of the individual jet rings.
At a given $s$, choose an $L$-basis adapted to $F_s$ from vectors
in $U_O$, clearing constant denominators. A scalar coefficient of
a vector of jump $a$ has order at least $-a$ there. Lift its jets
by polynomial interpolation, giving it order at least $H_t$ at
each other level $t$ by the Chinese remainder theorem on the
numerator. The resulting scalar times the chosen vector is an
allowed native section: take $\alpha_s=-a$, $\alpha_t=H_t$;
the lower jump bound gives $F_t^{-H_t}=D_c$. Any residual
multiple of $\Lambda$ is likewise a sum of allowed sections.
This argument works over $L$ itself by using the $P_s$ jet factors,
since the filtration at each level is the scalar extension of
$F_s$. Clearing the finitely many $p$-denominators gives the
single integer in $\mathrm{InNL}$. Integrality of the unscaled
numerator is already required in $\Xi$.

The exceptional quotient is the fixed module for $\gamma$ on
\[
      (\Xi_r+v\mathcal N)/(v\mathcal N),\qquad
                   \Xi_r=\Lambda^{-1}A\otimes S_r.
\]
It is finite free over $O$ of rank at most $(e+1)k$. First,
$v\mathcal N$ is $O$-saturated in $\Xi$: divisibility of the
numerator by $v$ can be tested at zero over $O$, and all remaining
conditions on the quotient by $v$ are tests in $L$-linear spaces.
Also $v\Lambda^2\Xi_r\subset v\mathcal N$, so distinguished
division gives finite generation. It remains to bound rational
invariants. After inverting $p$ the source $\Xi_r/v\Lambda^2\Xi_r$
splits into its $v$ factor and one family of jets per level $s$.
The quotient inherits an invariant filtration by Taylor degree on
each level, or equivalently one may filter successively by $v$
and the stable ideals $(P_s)$. At zero the contributing dimension
is at most $k$. After splitting coefficients, the degree-$i$ grade
at level $s$ is a quotient of a space of $k$-vectors on the tame
orbits of primitive $p^{s+1}$ roots, with action $u^i$ times
permutation in logarithmic coordinates, $u=1+p$. Equivalently
one takes the tame invariants in the degree-$i$ coefficient space,
where the tame substitution also acts by its $i$th power.
The $p^s$th power of $\gamma$ acts by $u^{i p^s}$, so only
$i=0$ can contribute. At $i=0$ the permutation on the tame orbits
is transitive. Its action has finite order in characteristic zero,
so a quotient of this grade has at most $k$ fixed dimensions.
Filtering the fixed subspace proves the stated rank bound.
A free basis has representatives in $\Xi_r$ satisfying
$(\gamma-1)f\in v\mathcal N$ because the fixed condition is
independent of adding $v\mathcal N$. This proves the exceptional
coefficient-generator assertion of $\mathrm{UnivCore}$ for this
actual arc. Its bound is $me(e+1)$, so the sum over all cuts is
at most $4dm$. This is a statement about coefficient generators;
realizing them by source or target classes is a further step.
% END U83M3

% BEGIN U83M4
\paragraph{U83M4. The shift polynomial and the cost still required of the scalar lifts.}
In the signed arc all types are up or down. Its shift set $J$ is
the union of the intervals $[-s_i-S,-s_i)$ on up letters and
$[s_i,s_i+S)$ on down letters, over its $e\le3$ letters.
In particular it does not contain $[-2G,G]$ from a collapsed thin
type. Apply K82P1--K82P2's polynomial construction with power
$p^e$ and the sum of the first $p^e$ powers. Its value at $1$
has normalized valuation
\[
 \kappa_c=e+\sum_{i\in J\setminus\{0\}}(e+1+v_p(i))
          \le 13S+3L_0+3 .
\]
Indeed every interval has at most $S$ integers with absolute value
at most $2h$. At level $p^a$ it meets at most $S/p^a+1$ nonzero
multiples, with none once $a>L_0$. Thus each interval contributes
at most $S/q+L_0$ to the valuation sum, while its other $e+1$
costs together are at most $12S$ over the arc; $3/q\le1$.
K82P1 gives $\gamma^{p^e}=u^{ip^e}$ on degree $i$ at every level,
which verifies annihilation of these shifted degrees on all branches.

For specificity, the residual implication for the general mixed
place can now be formulated directly in terms of finite classes.
Put $I$ equal to the native indicator of the $320$ band on the
word, $E_1=\lfloor40U+\eta I\rfloor$, and
$\kappa=\max_c\kappa_c$. Suppose the cut classes produced in
U83M2 can, after $p^{\sigma+\kappa_c+E_1}$, be lifted through
$F$ up to a fixed submodule generated by at most $me(e+1)$
integral finite target classes per cut. This is the remaining
cut-level comparison. Then $p^E$ times the entire quotient lies
in the span of at most $4dm$ such classes, with
\[
 E=R+3\sigma+\kappa+E_1\le55U+13S+\eta I .
\]
Indeed $3L_0+3\le3U$, and the floor and maxima only enlarge the
individual cut multipliers when brought to the common exponent.
All signs in their sum are harmless. Let $e'$ be the coarse
annihilation exponent from I82B1, at most $25h_B/p+3$. The
images of the exceptional classes are killed by that same power,
so additivity of normalized DVR length gives the upper bound
$(d+1)rE+4dm e'$. Here the free-plus-torsion presentation of
I82B1 covers all integral target classes; the onto lattice map
lifts their integral coboundaries. This fits the original row:
$U\le h_B/p^2+(w_P+\eta_3)/p+S_c+4+\mathrm{large}_2$,
$\log_p(2h)\le 3+\mathrm{large}_2$ as $h_B\le2b$, and
$\eta I\le\eta_3\,\mathrm{bandOf}(w)$ because the present good
reading witnesses the existential place indicator whenever $I=1$. Since
$\mathrm{large}_2\ge1$, the high term costs at most $55$ of its
coefficient $368$, the flat term at most $275$ of its coefficient
$2082$ (including $13S$), and $4 e'\le100(1+h_B/p+\log_p h_B)$
uses at most $100$ of its coefficient $22365$.

U83M3 supplies the rank on the constant-coordinate quotient.
The cut comparison above still requires transporting these
generators integrally into the physical module and lifting each
remaining one-coordinate section after multiplication by
$p^{E_1}$. On an arc of length two the source distinguishes the
ordering of $p(s_{c-1}+S+\eta)+\eta$ and $s_c$ in the terminal
up case, and compares $s_{c-1}$ with the local index plus $\eta+1$
in the terminal down case. On length three each initial sign
again gives two tests. The bound for coefficients at the cut
alone does not verify these intertwining and scalar-lifting tests.
% END U83M4

% BEGIN U83M5
\paragraph{U83M5. An unconditional mixed-place bound up to the larger exponential scale.}
Return to the actual coefficient test, and write $g$ for the
orbit maximum of M82R1, which can be smaller than $G+\eta$.
On every mixed orbit the native row bound already holds whenever
\[
                 p r g\le40m h_B .
\]
In fact M82R1 bounds the actual normalized cokernel by
$(d+1)r(32h/q^2+366g+14\log_p h)$. Spend its first term against
$368(d+1)rh_B/p^2$, since $p/q\le11/10$ and $h\le h_B$.
Spend the last against $2082(d+1)rF_w$ as in M82R2, where
$F_w\ge\mathrm{large}_2$. For the middle term use the third
summand of the frozen row interval, with
$E_p=1+h_B/p+\log_p h_B$. The even orbit has $d\ge2$, so
\[
 366(d+1)rg\le366\cdot\tfrac32\cdot40\,dmh_B/p
       =21960\,dmh_B/p\le22365\,dm E_p .
\]
This range requires no comparison from
U83M4. In particular every mixed reading with
\[
                   p\le \frac{\exp(x^{\epsilon/3})}{x^6}
\]
is covered uniformly on the late original family. Indeed
$r\le2x^5$ by A78P2, $g\le2b'$ by I82W3, and
$b'/b=x\exp(-x^{\epsilon/3})$ by I82W1. Therefore
$prg\le4b\le40m h_B$ throughout this range. More precisely
one can use the displayed orbit test on every good reading at
a place. When $f_p\ne\mathrm{cc}$, I82W5 excludes entirely
thick readings, and N82I1 supplies all the entirely nonthick
ones. G82I2 consequently proves the native global inequality at
these inert C or D indices up to the stated scale. The previously
proved test at larger primes in N82I2 continues to apply as well.
% END U83M5

% BEGIN U83M6
\paragraph{U83M6. Realizing the exceptional coefficient generators as finite target classes.}
Retain the integral constant basis of U83M3. This argument applies
to either tested lattice on one full signed arc. Let $i=c-e$,
write $b$ for its row of Wach basis vectors at $i$, $D_b=\varphi^e b$
for the translated row at $c$, and $U_b$ for the normalized
invariant-frame matrix at $v=0$. Its constant crystalline basis
at $c$ is $d_b=\varphi_{\rm cris}^e\mathrm{cmp}(b)$.
Put
\[
    f'=\varphi^{e-1}(v^M P_0^{2G})
       \sim v^M\prod_{s<e-1}P_s^{qM}P_{e-1}^{2G} .
\]
Since $qM>2h$ and $M\ge1$, $v\Lambda^2$ divides $f'$ at
all these levels, including the last one when $e=1$.
U83M2's monomial estimate applies to every entry of $U_b$,
whose positive-degree denominators satisfy the same recursion;
its constant term is integral. Apply $\varphi^{e-1}$ to the
integral representatives after the first Frobenius. This gives
an integral matrix $C$ equal to $p^\sigma\varphi^e(U_b\bmod v^M)$
in $(A/(f'))[1/p]$. For $g\in\Xi$ define the physical vector
\[
     V(g)=D_b C(\Lambda g)/\Lambda\in
              \Lambda^{-1}A\varphi^e L_i .
\]
This is $A$-linear. It represents $p^\sigma g$ in the precise
$\mathrm{RepVV}$ sense: choose one $p^k$ clearing the coefficients
of $U_b$ below degree $M$ to obtain its integral truncated columns,
and take $p^\sigma\Lambda g$ as the native numerator. The
congruence just proved makes the difference of the two resulting
numerators a multiple of $v\Lambda^2$ in the translated span,
after clearing any remaining $p$-denominator. This gives both
integers and the single modulus of the native representation.

At every primitive branch of level $s<e$, J82F2's substitution
identifies $D_b\varphi^e U_b=d_b$ in the full Taylor expansions.
Truncation at $M$ introduces error of order at least $qM$, since
the substituted tame parameter vanishes to order $q$ there. Thus
$V(g)$ differs in crystalline coordinates from $p^\sigma d_b g$
by an error of order at least $qM-H_s$ for $s<e-1$, and
$2G-H_s$ for $s=e-1$. Both are at least $H_s$. Every term of
such an error belongs to the true local filtered lattice since
all jumps there are at least $-H_s$. In particular
$V(\mathcal N)\subset L_c$ integrally. To see integrality here,
one first clears the factors $P_s$ by the two height inclusions
of U83M2. Divisibility in the free integral module can be tested
at all primitive roots of each distinguished factor, exactly as
in D82P1. The same argument applies to any vector in the
physical module $\Lambda^{-1}A\varphi^eL_i$.

Suppose $g\in\Xi$ and $(\gamma-1)g\in v\mathcal N$. Then
$(\gamma-1)V(g)/v$ belongs to the same physical module. Indeed
$\gamma$ preserves $(\Lambda)$, and in the basis $D_b$ the
numerator of $(\gamma-1)V(g)$ vanishes at $v=0$. The basis
action there is the identity by commutation with Frobenius,
while $\gamma(\Lambda)/\Lambda$ has value one there since
$\Lambda(0)\ne0$. At level $s$ the Taylor error of $V(g)$
and its generator translate still have order at least $H_s$;
$v$ is a Taylor unit. Division by $v$ therefore leaves the
same filtered test as $p^\sigma d_b(\gamma-1)g/v$, up to a
vector in the true local lattice. The preceding integral criterion
now gives $(\gamma-1)V(g)\in vL_c$. Thus $V(g)$ is a permitted
physical cut vector realizing $p^\sigma g$. This proves the
realization clause at the end of $\mathrm{UnivCore}$ on this arc.

Associate to $V(g)$ a cut class as follows. Apply U83M2's
truncation to $(\gamma-1)V(g)$, use
$p^\sigma\varphi V(g)-\tau$ as its first entry on $c+1$, and
solve for the second entry in $v^ML$ by the contracting
Frobenius difference. The generator difference of the first entry
lies in $v^ML$ by U83M2, and its poles were bounded above, so
H82A1 gives a finite class. The class depends $O$-linearly on
$g$ and vanishes if $g\in v\mathcal N$. In fact $V$ is linear
and maps that submodule into $vL_c$. For a difference in $vL_c$
the primitive is the vector itself modulo $v^M$ by uniqueness on
positive grades. Since $v^MP_0^{2G}\mid\varphi(v)^M$, the resulting
first entry belongs to $v^MP_0^{2G}A\varphi L_c\subset v^ML_{c+1}$.
The two entries then form a boundary by the contraction, as in
H82A2. The same argument with zero difference proves independence
of all truncation choices.

In particular representatives of the at most $me(e+1)$ basis
elements in U83M3 give actual finite target classes. These generate
all cut classes obtained in this way from its exceptional fixed
quotient, including its nontrivial integral lattice. The remaining
comparison in U83M4 must separate an arbitrary cut class into
these exceptional directions and the complementary one-coordinate
sections, and lift the latter through the specified source map.
% END U83M6
